\documentclass[10pt,a4paper]{article}
\usepackage{amsmath,amssymb,amsthm,graphicx}
\usepackage{xfrac}
\newenvironment{enumeratei}
{\bgroup\def\theenumi{\roman{enumi}}\def\theenumii{\arabic{enumii}}\begin{enumerate}}
{\end{enumerate}\egroup}

\hyphenation{addi-tion-al belong belongs denote exchan-ge exchan-ged exchan-ges}

\newcommand{\psfrac}[2]{\sfrac{(#1)}{#2}}

\newcommand\mto{\mathchoice{\longmapsto}{\mapsto}{\mapsto}{\mapsto}}
\let\oldbigsqcup\bigsqcup
\renewcommand{\bigsqcup}{\mathop{\textstyle\oldbigsqcup}\displaylimits}

\let\epsilon\varepsilon

\DeclareMathOperator{\Ess}{Ess}
\DeclareMathOperator{\Leb}{Leb}

\usepackage{tikz}
\usetikzlibrary{plotmarks}

\theoremstyle{plain}
\newtheorem{theorem}{Theorem}[section]
\newtheorem{lemma}[theorem]{Lemma}
\newtheorem{proposition}[theorem]{Proposition}
\newtheorem{corollary}[theorem]{Corollary}
\newtheorem{claim}{Claim}

\theoremstyle{definition}
\newtheorem{remark}[theorem]{Remark}
\newtheorem{example}[theorem]{Example}

\newcommand{\A}{\mathcal{A}}
\newcommand{\B}{\mathcal{B}}
\newcommand{\I}{\mathcal{I}}

\newcommand{\Z}{\mathbb{Z}}
\newcommand{\N}{\mathbb{N}}
\newcommand{\R}{\mathbb{R}}
\title{Backward iterates of the reflection centre are given by the conjugacy with the inverse map}
\author{Przemyslaw Berk, Frank Trujillo}
\date{doi:10.5802/jep.302; CC BY4.0}
\begin{document}
\maketitle
\noindent\textit{Editorial note.} The counted statement is the author's labelled lemma environment at source0.tex lines 530-545; its proof environment opens at line 547, the line after the blank separator, so the pairing is by adjacency and the retained proof range 547-549 is the whole environment.\par\medskip
\noindent\textit{Editorial note.} The numbered identity eq:generalized\_conjugacy (source lines 469-472, a multline environment) is retained as context because the counted proof applies it at source line 548; it is the only outside label of the retained proof.\par\medskip
\noindent\textit{Editorial note.} This source is new to the material-authoring work and its preamble is built from the author's own macro and theorem block (source lines 8-45) rather than from a prepared adapter.\par\medskip
\noindent\textit{Editorial note.} This article carries no label-number dump and the retained counter values are left empty for the whole source, so the printed numbering of the card is the wrapper's default; the identity of the counted item is its author label lem:inverse\_iterates.\par\medskip
\noindent\textit{Editorial scope.} One original result of the pinned publisher article, counted once: the formula for the backward iterates of a reflection centre of a symmetric interval exchange transformation (author label lem:inverse\_iterates).\par\medskip
\noindent\textit{Allowed standard prerequisites (editorial).}\begin{itemize}
\item Symmetric interval exchange transformation. An IET is symmetric when the reflection of the interval conjugates it to its inverse off the discontinuities.. It is the hypothesis and the source of the conjugacy identity.
\item Reflection centres. c\_alpha denotes the centre of the subinterval indexed by alpha, and c\_\{1/2\} is the centre of the whole interval.. They are the points whose backward iterates are computed.
\item Generalised conjugacy identity. The reflection applied to the n-th iterate equals the minus-n-th iterate of the reflection, off the relevant discontinuities.. It is the identity applied in the proof.
\end{itemize}\medskip
\begin{multline}
\label{eq:generalized_conjugacy}
\mathcal{I}_I\circ T^n(x)=T^{-n}\circ\mathcal{I}_I(x) \quad \text{if } x \neq T^{-i}(\partial I_\alpha) \text{ for any } \alpha \in \A \text{ and}\\ \text{any $i$ such that } 0 \leq i < n.
\end{multline}

\begin{lemma}
\label{lem:inverse_iterates}
Let $T: I \to I$ be a symmetric IET and $\alpha \in \A \cup \{\sfrac{1}{2}\}$. If $m \geq 1$ is such that $\{c_\alpha, T(c_\alpha), \ldots, T^{m - 1}(c_\alpha)\} \cap \{\partial I_\beta\}_{\beta \in \A} = \emptyset$, then\vspace*{-3pt}
\begin{equation}
\label{eq:inverse_iterates}
T^{-m}(c_\alpha) =
{T^{-1} \circ \I_I} \big( T^{m - \delta_{1/2}(\alpha)}(c_\alpha) \big),
\end{equation}
where\vspace*{-3pt}
\begin{equation}
\label{eq:delta_1/2}
\delta_{1/2}(\beta) = \left\{\begin{array}{ll} 1 & \text{ if } \quad \beta = \sfrac{1}{2}, \\
0 & \text{ otherwise}.
\end{array} \right.
\end{equation}
\end{lemma}

\begin{proof}
This result follows by directly applying \eqref{eq:generalized_conjugacy} and noticing that $\I_I(c_{1/2}) = c_{1/2}$ and that $\I_I(c_\alpha) = T(c_\alpha)$, for any $\alpha \in \A$.
\end{proof}

\end{document}
